
\documentclass{article}
\usepackage{colortbl}
\usepackage{makecell}
\usepackage{multirow}
\usepackage{supertabular}

\begin{document}

\newcounter{utterance}

\centering \large Interaction Transcript for game `taboo', experiment `low\_en', episode 2 with reka{-}flash{-}3.1{-}t0.0.
\vspace{24pt}

{ \footnotesize  \setcounter{utterance}{1}
\setlength{\tabcolsep}{0pt}
\begin{supertabular}{c@{$\;$}|p{.15\linewidth}@{}p{.15\linewidth}p{.15\linewidth}p{.15\linewidth}p{.15\linewidth}p{.15\linewidth}}
    \# & $\;$A & \multicolumn{4}{c}{Game Master} & $\;\:$B\\
    \hline

    \theutterance \stepcounter{utterance}  
    & & \multicolumn{4}{p{0.6\linewidth}}{
        \cellcolor[rgb]{0.9,0.9,0.9}{
            \makecell[{{p{\linewidth}}}]{
                \texttt{\tiny{[P1$\langle$GM]}}
                \texttt{You are playing a collaborative word guessing game in which you have to describe a target word for another player to guess.} \\
\\ 
\texttt{Rules:} \\
\texttt{(a) You have to reply in the form: CLUE: <some text>. Guesses from the other player will start with GUESS.} \\
\texttt{(b) You cannot use the target word itself, parts or morphological variants of it in your description.} \\
\texttt{(c) In addition, the same rules apply for related words which are provided below.} \\
\\ 
\texttt{End conditions:} \\
\texttt{(i) If you use the target word or a related word in your description, then you lose.} \\
\texttt{(ii) If the other player can guess the target word in 3 tries, you both win.} \\
\\ 
\texttt{Let us start.} \\
\\ 
\texttt{This is the target word that you need to describe and that the other player needs to guess:} \\
\\ 
\texttt{competent} \\
\\ 
\texttt{Related words are:} \\
\\ 
\texttt{{-} efficient} \\
\texttt{{-} skilled} \\
\texttt{{-} capable} \\
\\ 
\texttt{Important: You are under time pressure, give short descriptions that are to the point!} \\
            }
        }
    }
    & \\ \\

    \theutterance \stepcounter{utterance}  
    & \multicolumn{4}{p{0.6\linewidth}}{
        \cellcolor[rgb]{0.9,0.9,0.9}{
            \makecell[{{p{\linewidth}}}]{
                \texttt{\tiny{[P1$\rangle$GM]}}
                \texttt{<reasoning>} \\
\texttt{I need to come up with a clue for the word "competent" without using the word itself or any related words like efficient, skilled, or capable. Let me think about what "competent" means. It means having the ability, knowledge, or skill to do something successfully. So maybe something about being good at something, but I can't use "skilled" or "capable."} \\
\\ 
\texttt{Hmm, maybe use a synonym in a different way. Wait, the rules say I can't use morphological variants. So words like "competence" or "competently" are out. Let me think of other attributes. Competent people are reliable, maybe? Or maybe they can handle tasks well. How about something like "adequate for the job"? But "adequate" might be too vague. Or "effective in performance"? But "effective" is related to efficient, which is a related word. Wait, the related words are efficient, skilled, capable. So I can't use those either. So I need to avoid those terms.} \\
\\ 
\texttt{What else defines competent? Maybe someone who can be trusted to do a job properly. How about "reliable and effective worker"? Wait, "effective" is related to efficient, which is a related word. So that's not allowed. Maybe "adequate performer"? Still not sure. Let me think of a different angle. Competent people can solve problems, make decisions. Maybe "problem{-}solver with good judgment"? But "good judgment" might be okay. Wait, the clue has to be short. Maybe "adequately skilled"? No, "skilled" is a related word. Hmm.} \\
\\ 
\texttt{Wait, maybe use a metaphor. Like "someone who can handle the job without issues." But "handle" is okay? Not sure. Alternatively, think of professions. "A professional who meets expectations." But "professional" might be okay. Or "someone who gets the job done right." "Gets the job done" is a common phrase. "Right" implies doing it correctly. Maybe that works. Let me check if any of the related words are in there. "Gets the job done right" – "right" isn't a related word. "Job done" is more about efficiency maybe? Not sure. But the related words are efficient, skilled, capable. So as long as I don't use those exact words, it's okay. "Gets the job done right" –} \\
            }
        }
    }
    & & \\ \\

    \theutterance \stepcounter{utterance}  
    & & & \multicolumn{2}{p{0.3\linewidth}}{
        \cellcolor[rgb]{0.9,0.9,0.9}{
            \makecell[{{p{\linewidth}}}]{
                \texttt{\tiny{[GM$|$GM]}}
                \texttt{abort game} \\
            }
        }
    }
    & & \\ \\

\end{supertabular}
}

\end{document}
